\subsection{Schwartz functions}

We denote \(\mathcal{S}(\mathbf{R}^{n})\) the space of infinitely differentiable functions \(\varphi\) on \(\mathbf{R}^{n}\) with complex values, such that, for all \(\alpha\) and \(\beta\) in \(\mathbf{N}^{n}\) , the function \(\mathrm{X}^{\beta}\partial^{\alpha}\varphi\) is bounded on \(\mathbf{R}^{n}\) . We equip \(\mathcal{S}(\mathbf{R}^{n})\) with the locally convex topology defined by the countable family of semi-norms \((q_{\alpha,\beta})_{(\alpha,\beta)\in\mathbf{N}^{n}\times\mathbf{N}^{n}}\) , where \(q_{\alpha,\beta}\) is defined by

\[
q _ {\alpha , \beta} (\varphi) = \sup _ {x \in \mathbf {R} ^ {n}} | x ^ {\beta} \partial^ {\alpha} \varphi (x) | = \| \mathrm{X} ^ {\beta} \partial^ {\alpha} \varphi \| _ {\infty}
\]

for \(\varphi\in\mathcal{S}(\mathbf{R}^{n})\) . This topology is Hausdorff. It is also defined by the semi-norms \(\widetilde{q}_{\alpha,k}\) defined by

\[
\widetilde {q} _ {\alpha , k} (\varphi) = \sup _ {x \in \mathbf {R} ^ {n}} \| x \| ^ {k} | (\partial^ {\alpha} \varphi) (x) |.
\]

for all \(\varphi\in\mathcal{S}(\mathbf{R}^{n})\) , where \(k\in\mathbf{N}\) and \(\alpha\in\mathbf{N}^{n}\) .

We say that \(\mathcal{S}(\mathbf{R}^{n})\) is the Schwartz space or the space of Schwartz functions on \(\mathbf{R}^{n}\) .

Note. --- Let \(\varphi\in\mathcal{S}(\mathbf{R}^{n})\) . For every \(k\in\mathbf{N}\) , we have

\[
\lim _ {\| x \| \to + \infty} \| x \| ^ {k} \varphi (x) = 0,
\]

since the function \(x \mapsto \|x\|^{k+1} \varphi(x)\) is bounded.

Example. --- The function \(\gamma_{n}\) defined on \(\mathbf{R}^{n}\) by \(\gamma_{n}(x)=\exp(-\|x\|^{2})\) belongs to \(\mathcal{S}(\mathbf{R}^{n})\) . Indeed, one proves by induction on k that, for every integer \(k\in \mathbf{N}\) , there exists a polynomial \(P_{k}\in\mathbf{R}[X]\) such that \(\partial_{k}\gamma_{1}=P_{k}\gamma_{1}\) .

For all \(\alpha = (\alpha_{i})\in \mathbf{N}^{n}\) and \(\beta = (\beta_{i})\in \mathbf{N}^{n}\) , and every \(x = (x_{i})\in \mathbf{R}^{n}\) , it then follows that

\[
| (\mathrm{X} ^ {\beta} \partial^ {\alpha} \gamma_ {n}) (x) | = \prod_ {i = 1} ^ {n} | x _ {i} | ^ {\beta_ {i}} | \mathrm{P} _ {\alpha_ {i}} (x _ {i}) | \gamma_ {1} (x _ {i}),
\]

which is a bounded quantity as x varies in \(\mathbf{R}^{n}\) .

Let \(\alpha \in \mathbf{N}^n\) and \(\beta \in \mathbf{N}^n\) . If \(\varphi \in \mathcal{S}(\mathbf{R}^n)\) , then \(\mathrm{X}^{\beta}\partial^{\alpha}(\varphi)\) is still a Schwartz function; the map \(\varphi \mapsto \mathrm{X}^{\beta}\partial^{\alpha}\varphi\) thus defined from \(\mathcal{S}(\mathbf{R}^n)\) into itself is continuous.

The space \(\mathcal{S}(\mathbf{R}^{n})\) is a topological algebra. More precisely, if \(\varphi_{1}\) and \(\varphi_{2}\) belong to \(\mathcal{S}(\mathbf{R}^{n})\) , then \(\varphi_{1}\varphi_{2}\) is a Schwartz function such that

\[
q _ {\alpha , \beta} (\varphi_ {1} \varphi_ {2}) \leqslant \sum_ {0 \leqslant \gamma \leqslant \alpha} \binom{\alpha}{\gamma} q _ {\gamma , \beta} (\varphi_ {1}) q _ {\alpha - \gamma , 0} (\varphi_ {2})\tag{2}
\]

for all \(\alpha\) and \(\beta \in \mathbf{N}^n\) (prop. 1 of IV, p. 196).

The canonical inclusion of \(\mathcal{S}(\mathbf{R}^{n})\) into the space \(\mathcal{C}^{\infty}(\mathbf{R}^{n})\) , equipped with the topology described in §3 of IV, p. 200, is continuous, since

\[
\sup _ {x \in K} | \partial^ {\alpha} \varphi (x) | \leqslant q _ {\alpha , 0} (\varphi)
\]

for every compact subset K of \(\mathbf{R}^{n}\) , all \(\alpha \in \mathbf{N}^{n}\) and every Schwartz function \(\varphi\) .

Lemma 7. --- Let \(k \in \mathbf{N}\) and \(\alpha \in \mathbf{N}^n\) . For every function \(\varphi \in \mathcal{S}(\mathbf{R}^n)\) and every real number \(T > 0\), we have

\[
\widetilde {q} _ {\alpha , k} (\varphi) \leqslant \mathrm{T} ^ {k} \sup _ {\| x \| \leqslant \mathrm{T}} | \partial^ {\alpha} \varphi (x) | + \frac {1}{\mathrm{T}} \widetilde {q} _ {\alpha , k + 1} (\varphi).\tag{3}
\]

Indeed, one has

\[
\begin{array}{c} \widetilde {q} _ {\alpha , k} (\varphi) \leqslant \sup _ {\| x \| \leqslant T} \| x \| ^ {k} | \partial^ {\alpha} \varphi (x) | + \sup _ {\| x \| > T} \| x \| ^ {k} | \partial^ {\alpha} \varphi (x) | \\ \leqslant T ^ {k} \sup _ {\| x \| \leqslant T} | \partial^ {\alpha} \varphi (x) | + \frac {1}{T} \sup _ {\| x \| > T} \| x \| ^ {k + 1} | \partial^ {\alpha} \varphi (x) |. \end{array}
\]

PROPOSITION 10. --- Let B be a bounded subset of \(\mathcal{S}(\mathbf{R}^{n})\) . The topology induced on B by \(\mathcal{S}(\mathbf{R}^{n})\) coincides with the topology induced by \(\mathcal{C}^{\infty}(\mathbf{R}^{n})\) .

Since the inclusion of \(\mathcal{S}(\mathbf{R}^{n})\) into \(\mathcal{C}^{\infty}(\mathbf{R}^{n})\) is continuous, it suffices to show that, for every open subset V of \(\mathcal{S}(\mathbf{R}^{n})\) , the intersection \(V \cap B\) is open in B for the topology induced by \(\mathcal{C}^{\infty}(\mathbf{R}^{n})\) .

Let V be an open subset of \(\mathcal{S}(\mathbf{R}^{n})\) . Let \(\varphi_{0} \in \mathrm{V} \cap \mathrm{B}\) . There exist a finite set I, a family \((\alpha_{i}, k_{i})_{i \in \mathrm{I}} \in (\mathbf{N}^{n} \times \mathbf{N})^{\mathrm{I}}\) and a real number \(\varepsilon > 0\) such that V contains the set of \(\varphi \in \mathcal{S}(\mathbf{R}^{n})\) satisfying

\[
\sup _ {i \in I} \widetilde {q} _ {\alpha_ {i}, k _ {i}} (\varphi - \varphi_ {0}) \leqslant \varepsilon .
\]

Since B is bounded in \(\mathcal{S}(\mathbf{R}^{n})\) , there exists \(M > 0\) such that the seminorms \(\widetilde{q}_{\alpha_{i},k_{i} + 1}\) for \(i\in I\) are bounded by \(M\) on \(B\) . Let \(\delta >0\) and \(T > 0\) be real numbers. By inequality (3), as soon as \(\varphi \in B\) satisfies the bound

\[
\sup _ {i \in I} \sup _ {\| x \| \leqslant T} | \partial^ {\alpha_ {i}} (\varphi - \varphi_ {0}) | \leqslant \delta ,\tag{4}
\]

we have

\[
\sup _ {i \in I} \widetilde {q} _ {\alpha_ {i}, k _ {i}} (\varphi - \varphi_ {0}) \leqslant \delta T ^ {k} + \frac {2 M}{T}.
\]

Set \(T = \frac{4M}{\varepsilon}\) , then \(\delta = \frac{\varepsilon}{2T^{k}}\) . One observes that \(V \cap B\) contains the neighborhood of \(\varphi_{0}\) in B for the topology induced by \(\mathcal{C}^{\infty}(\mathbf{R}^{n})\) which is defined by (4). This concludes the proof.

COROLLARY. --- Let \((\varphi_{m})_{m\in\mathbf{N}}\) be a bounded sequence in \(\mathcal{S}(\mathbf{R}^{n})\) . For every function \(\varphi\in\mathcal{S}(\mathbf{R}^{n})\) , the following assertions are equivalent:

\begin{enumerate}
\def\labelenumi{\alph{enumi})}
\item
  The sequence \((\varphi_{m})\) converges to \(\varphi\) in \(\mathcal{S}(\mathbf{R}^n)\) ;
\item
  The sequence \((\varphi_{m})\) converges to \(\varphi\) in \(\mathcal{C}^{\infty}(\mathbf{R}^{n})\) .
\end{enumerate}

Remark. --- A sequence \((\varphi_{m})\) in \(\mathcal{C}^{\infty}(\mathbf{R}^{n})\) converges if and only if, for every \(\alpha\in \mathbf{N}^{n}\) the sequence \((\partial^{\alpha}\varphi_{m})\) converges to a function \(\varphi^{(\alpha)}\) in \(\mathcal{C}(\mathbf{R}^{n})\) endowed with the topology of compact convergence. We then have \(\varphi^{(\alpha)}=\partial^{\alpha}\varphi\) and \((\varphi_{m})\) converges to \(\varphi^{(0)}\) .

Indeed, the condition is necessary. Conversely, if the sequences \((\partial^{\alpha}\varphi_{m})\) converge to functions \(\varphi^{(\alpha)}\) for all \(\alpha\in \mathbf{N}^{n}\) , it then follows from FVR, II, p. 2, th. 1, that \(\varphi^{(\alpha)}=\partial^{\alpha}\varphi^{(0)}\) , which means that the sequence \((\varphi_{m})\) converges to \(\varphi^{(0)}\) in \(\mathcal{C}^{\infty}(\mathbf{R}^{n})\) .

PROPOSITION 11. --- The space \(\mathcal{S}(\mathbf{R}^{n})\) is a Fréchet space and a Montel space.

As the space \(\mathcal{C}^{\infty}(\mathbf{R}^{n})\) is complete (EVT, III, p. 9, example \(b\) ), the corollary of Proposition 10 implies that every Cauchy sequence in \(\mathcal{S}(\mathbf{R}^{n})\) converges in \(\mathcal{S}(\mathbf{R}^{n})\) since it is bounded and converges in \(\mathcal{C}^{\infty}(\mathbf{R}^{n})\) .

The space \(\mathcal{S}(\mathbf{R}^n)\) is therefore a Fréchet space; in particular, it is barrelled (EVT, III, p. 25, cor. of prop. 2). Let B be a bounded subset of \(\mathcal{S}(\mathbf{R}^n)\) and \((\varphi_m)_{m \in \mathbf{N}}\) a sequence with values in B. Since \(\mathcal{C}^\infty(\mathbf{R}^n)\) is a Montel space (EVT, IV, p. 18, example (4)), there exists a subsequence of \((\varphi_m)_{m \in \mathbf{N}}\) which converges in \(\mathcal{C}^\infty(\mathbf{R}^n)\) , hence in \(\mathcal{S}(\mathbf{R}^n)\) (proposition 10). Hence B is relatively compact in \(\mathcal{S}(\mathbf{R}^n)\) . It follows that \(\mathcal{S}(\mathbf{R}^n)\) is a Montel space.

